\subsection{Poincaré groupoid of a quotient space}

Let X be a topological space equipped with a continuous action of a discrete group G; set Y = X/G and denote by \(f: X \to Y\) the canonical map. Denote by \(\lvert G\rvert\) the groupoid \(\varpi(G)\); its set of vertices is G; for \(g, g' \in G\) , there exists a unique arrow from g to \(g'\) if \(g' = g\) and none otherwise. By passing to fundamental groupoids, the operation \(m: G \times X \to X\) induces a morphism of groupoids \(\varpi(m): \lvert G\rvert \times \varpi(X) \to \varpi(X)\) . Let \(\varpi(X)/G\) be the coequalizer of the two morphisms of groupoids \(\varpi(m)\) and \(\varpi(\mathrm{pr}_{2})\) of \(\lvert G\rvert \times \varpi(X)\) into \(\varpi(X)\) ; let us denote by \(\beta: \varpi(X) \to \varpi(X)/G\) the canonical morphism of groupoids. We have \(f \circ m = f \circ pr_{2}\) , therefore \(\varpi(f) \circ \varpi(m) = \varpi(f) \circ \varpi(\mathrm{pr}_{2})\) . By the universal property of coequalizers, there therefore exists a unique morphism of groupoids \(\varpi''(f): \varpi(X)/G \to \varpi(Y)\) such that one has \(\varpi(f) = \varpi''(f) \circ \beta\) .

THEOREM 3. --- Let X be a deloopable topological space and let G be a discrete group acting properly on X; let \(f\colon X \to X/G\) be the canonical surjection. The canonical groupoid morphism \(\varpi''(f)\colon \varpi(X)/G \to \varpi(X/G)\) introduced above is an isomorphism.

Denote by \(\operatorname{Coeg}(f)\) the coequalizer of the two groupoid morphisms from \(\varpi(\mathrm{X} \times_{\mathrm{Y}} \mathrm{X})\) into \(\varpi(\mathrm{X})\) induced by the projections \(pr_{1}\) and \(pr_{2}\) ; let \(\gamma: \varpi(\mathrm{X}) \to \operatorname{Coeg}(f)\) be the canonical groupoid morphism. The image of the map \((m, pr_{2})\) : \(G \times X \to X \times X\) being the subspace \(X \times_{Y} X\) of \(X \times X\) , the two groupoid morphisms \(\gamma \circ \varpi(m)\) and \(\gamma \circ \varpi(pr_{2})\) of \(\lvert G\rvert \times \varpi(X)\) into \(\operatorname{Coeg}(f)\) are equal. By the universal property of \(\varpi(\mathrm{X})/\mathrm{G}\) , there exists a unique morphism of groupoids \(\alpha: \varpi(\mathrm{X})/\mathrm{G} \to \operatorname{Coeg}(f)\) such that \(\gamma = \alpha \circ \beta\) .

We also denote by \(\varpi'(f)\) the unique groupoid morphism from \(\operatorname{Coeg}(f)\) into \(\varpi(Y)\) such that \(\varpi(f) = \varpi'(f) \circ \gamma\) . By IV, p.~399, example 2, the hypotheses of th. 1 of IV, p.~398 are satisfied and the morphism \(\varpi'(f)\) is an isomorphism. Since

\[
\varpi (f) = \varpi^ {\prime} (f) \circ \gamma = \varpi^ {\prime} (f) \circ \alpha \circ \beta = \varpi^ {\prime \prime} (f) \circ \beta ,
\]

we have \(\varpi''(f) = \varpi'(f) \circ \alpha\) . Thus, to prove Theorem 3, it suffices to show that the morphism \(\alpha\) is an isomorphism.

Lemma 4. --- For every path \(c = (c_1, c_2)\) in \(X \times_Y X\) , we have \(\beta([c_1]) = \beta([c_2])\) .

Let c be such a path.

Let \(x \in X\) , let \(K_x\) be its stabilizer in G. By TG, III, p.~32, prop. 8, there exists an open neighborhood \(U_x\) of \(x\) in X such that \(K_x \cdot U_x = U_x\),

\(g \cdot U_x \cap U_x = \varnothing\) for all \(g \in G - K_x\) and such that the map f induces a homeomorphism from \(U_x/K_x\) onto an open neighborhood \(V_x\) of \(f(x)\) in Y. Since X is deloopable and the restriction to \(U_x\) of the map f is open and closed, one may moreover assume that \(U_x\) is connected and that the image of the canonical homomorphism \(\pi_1(U_x, x) \to \pi_1(X, x)\) is reduced to the identity element. The open sets \((V_x)_{x \in X}\) thus constructed form an open covering of Y. By Lemma 4 of III, p.~272, applied to the compact space I and the open sets \((f \circ c_1)^{-1}(V_x)\) for \(x \in X\) , there exists an integer \(n \geqslant 1\) such that for every \(i \in \{1, \ldots, n\}\) , there exists a point \(x_i\) in X such that \(c_1([i-1/n, i/n])\) is contained in \(f(V_{x_i})\) . Since \(f \circ c_1 = f \circ c_2\) , \(c_2([i-1/n, i/n])\) is also contained in \(f(V_{x_i})\) .

For \(j = 1\) or 2 and for \(i \in \{1, \ldots, n\}\) , denote by \(c_{j,i}\) the path in X defined by \(t \mapsto c_j(\frac{i+t-1}{n})\) ; we have \(c_j = c_{j,1} * \cdots * c_{j,n}\) (III, p.~291, remark 1). Consequently, to show that \(\beta([c_1]) = \beta([c_2])\) , it suffices to show that \(\beta([c_{1,i}]) = \beta([c_{2,i}])\) for every integer \(i \in \{1, \ldots, n\}\) . Replacing the path \((c_1, c_2)\) by the path \((c_{1,i}, c_{2,i})\) , we may thus assume that there exists \(x \in X\) such that \((f \circ c_1)([0,1]) \subset V_x\) .

The inverse image of \(V_{x}\) under f is the disjoint union of the connected components \(g \cdot U_{x}\) where g runs over a system of representatives in G of \(G/K_{x}\) . For i = 1 or 2, let \(g_{i}\) be an element of G such that the point \(x_{i} = g_{i} \cdot c_{i}(0)\) belongs to \(U_{x}\) . The image of the path \(g_{i} \cdot c_{i}\) is then contained in \(U_{x}\) . By definition of the groupoid \(\varpi(\mathrm{X})/\mathrm{G}\) , we have \(\beta([g_{i} \cdot c_{i}]) = \beta([c_{i}])\) which allows us to assume that the images of the paths \(c_{1}\) and \(c_{2}\) are contained in \(U_{x}\) and that \(g_{1} = g_{2} = e\) .

For \(s=0\) or 1, let \(d_{s}\) be a path in \(\mathrm{U}_{x}\) connecting \(x\) to \(c_{1}(s)\) and let \(g_{s}\) be an element of \(\mathrm{K}_{x}\) such that \(g_{s}\cdot c_{2}(s)=c_{1}(s)\) . The paths \(d_{0}*c_{1}*\overline{d_{1}}\) and \(d_{0}*(g_{0}\cdot c_{2})*(g_{0}g_{1}^{-1}\cdot\overline{d_{1}})\) are loops at \(x\) on \(\mathrm{U}_{x}\) . They are therefore strictly homotopic in X to the constant loop at \(x\) , because the image of the canonical homomorphism \(\pi_{1}(\mathrm{U}_{x},x)\to\pi_{1}(\mathrm{X},x)\) is reduced to the identity element. Their classes in particular have the same image under the groupoid morphism \(\beta\) , whence

\[
\beta ([ d _ {0} ]) \beta ([ c _ {1} ]) \beta ([ d _ {1} ]) ^ {- 1} = \beta ([ d _ {0} ]) \beta ([ g _ {0} \cdot c _ {2} ]) \beta ([ g _ {0} g _ {1} ^ {- 1} \cdot d _ {1} ]) ^ {- 1}.
\]

Since \(\beta([g \cdot c]) = \beta([c])\) for every element \(g \in G\) and for every path \(c\) in \(X\) , the equality \(\beta([c_1]) = \beta([c_2])\) follows, as was to be proved.

By the lemma, the two groupoid morphisms \(\beta \circ \varpi(\mathrm{pr}_{1})\) and \(\beta \circ \varpi(\mathrm{pr}_{2})\) from \(\varpi(\mathrm{X} \times_{\mathrm{Y}} \mathrm{X})\) into \(\operatorname{Coeg}(f)\) are equal. By the universal property of the coequalizer, there therefore exists a unique groupoid morphism \(\alpha'\) : \(\operatorname{Coeg}(f) \to \varpi(\mathrm{X})/\mathrm{G}\) such that \(\beta = \alpha' \circ \gamma\) . The morphism \(\alpha' \circ \alpha\) is the unique groupoid morphism \(\varphi\) of \(\varpi(\mathrm{X})/\mathrm{G}\) into itself such that \(\varphi \circ \beta = \beta\) ; one therefore has \(\alpha' \circ \alpha = \operatorname{Id}_{\varpi(\mathrm{X})/\mathrm{G}}\) . Likewise, \(\alpha \circ \alpha' = \operatorname{Id}_{\operatorname{Coeg}(f)}\) . Consequently, \(\alpha\) is an isomorphism.

